\section{Data Analysis I: Market Demand and Pay (RQ1)}\label{sec:rq1}

\subsection{Demand}

Within the 5,138 data-role postings, the most requested skills (with Wilson 95\% intervals
\citep{wilson1927}) are data analytics (18.3\% [17.3, 19.4]), SQL (16.5\% [15.5, 17.5]), SAS (13.4\%),
Python (12.6\%), business analysis (12.5\%), machine learning (12.0\%) and R (10.8\%)
(Figure~\ref{fig:fig_demand.png}). The market expects a \emph{toolkit}: one statistical language
(SAS, R or Python) plus SQL.

\fig[0.78\linewidth]{fig_demand.png}{Most demanded skills in data-role postings (Wilson 95\% CIs).}

\subsection{Pay, after accounting for experience}

Raw salary comparisons are misleading because skills and seniority go together. We therefore fit an
ordinal logistic regression of the six salary bands on indicators for the 30 most common skills,
minimum experience, role family and a multi-city flag. The odds ratio (OR) of a skill is the
multiplicative change in the odds of being in a \emph{higher} band when a posting lists it, at the
same experience and role family.

\textbf{Experience dominates:} each extra year of required experience multiplies the odds of a
higher band by 1.66 (all postings) and 1.61 (data roles). Beyond experience, 13 skills remain
significant at FDR $q<0.05$ in data roles (Table~\ref{tab:pay}, Figure~\ref{fig:fig_pay_data_roles.png}).

\begin{table}[H]\centering\small
\caption{Skills associated with salary band in data-role postings (ordinal logit, adjusted for
experience, role family and multi-city; FDR $q<0.05$).}\label{tab:pay}
\begin{tabular}{@{}l c c l c@{}}
\toprule
Higher-pay skill & OR [95\% CI] & & Lower-pay skill & OR [95\% CI]\\
\midrule
Product management & 2.95 [2.13, 4.07] & & MIS & 0.25 [0.19, 0.34]\\
Natural language processing & 1.74 [1.25, 2.41] & & Communication & 0.70 [0.55, 0.89]\\
Predictive modelling & 1.67 [1.28, 2.18] & & & \\
Marketing analytics & 1.64 [1.20, 2.26] & & & \\
Java & 1.51 [1.20, 1.90] & & & \\
Statistics & 1.43 [1.09, 1.88] & & & \\
SAS & 1.42 [1.20, 1.68] & & & \\
Business intelligence / Tableau & 1.41 / 1.41 & & & \\
Data mining & 1.40 [1.07, 1.83] & & & \\
Data analytics & 1.25 [1.09, 1.42] & & & \\
\bottomrule
\end{tabular}
\end{table}

\fig[0.82\linewidth]{fig_pay_data_roles.png}{Adjusted pay association of skills in data roles
(only skills with FDR $q<0.05$).}

\textbf{Interpretation.} Advanced analytics (predictive modelling, NLP, statistics, data mining)
and core tools (SAS, SQL, Java) are linked to higher bands \emph{beyond what experience explains}.
Postings emphasising MIS reporting or generic ``communication'' are lower-paid support roles; this
does not mean communication is unvaluable, as RQ2 shows.

\subsection{Skill bundles}

Association rules show that the big-data stack is requested as a bundle: hive $\rightarrow$ Hadoop
(confidence 75\%, lift 13.1), Scala $\rightarrow$ Spark (72\%, lift 15.2), hive $\rightarrow$ Spark
(68\%, lift 14.3); deep learning and NLP co-occur with lift 13.6 (Figure~\ref{fig:fig_bundles.png}).
These skills should be taught together, which is why the RQ4 catalogue contains a single Big-Data
course.

\fig[0.85\linewidth]{fig_bundles.png}{Strongest skill co-occurrence rules in data roles.}

\subsection{Company-level view (DataScience Jobs)}

Average salary rises with minimum experience (Spearman $\rho = 0.63$, $p<10^{-100}$). Openings are
spread across many employers (Herfindahl index 0.021), but the ten largest (TCS, Accenture,
Cognizant, Wipro, IBM, Genpact, Capgemini, LTI, Tech Mahindra, HCL) offer 35\% of all openings. The
job-weighted average salary (9.7 lakh) is below the per-row median (11.9 lakh): \emph{most openings
are at large IT-services firms that pay below the typical company} (Figure~\ref{fig:fig_ds_jobs.png}).

\fig{fig_ds_jobs.png}{DataScience Jobs: salary vs experience (bubble size = openings) and the ten
largest employers.}

\subsection{Triangulation: the five JDS skill areas in the market}

To connect RQ1 with RQ2, each of the five JDS skill areas was measured in the job market
(Table~\ref{tab:tri}, Figure~\ref{fig:fig_triangulation.png}).

\begin{table}[H]\centering\small
\caption{The five JDS skill areas: market demand and pay (RQ1, data roles) vs junior salary-hike
link (RQ2).}\label{tab:tri}
\begin{tabular}{@{}l c c c@{}}
\toprule
Skill area & Demand (\% postings) & Market pay OR [95\% CI] & RQ2 hike OR per SD [95\% CI]\\
\midrule
Coding & \textbf{40.7} & 1.50 [1.34, 1.68] & 1.67 [1.10, 2.79]\\
Maths \& stats & 15.9 & \textbf{1.56} [1.35, 1.80] & \textbf{3.01} [1.99, 4.60]\\
AI \& ML & 16.0 & 1.44 [1.23, 1.69] & 2.00 [1.23, 3.23]\\
Big data & 10.6 & 1.33 [1.10, 1.62] & 1.76 [1.18, 2.67]$^\dagger$\\
Dashboards \& storytelling & 11.9 & 1.14 [0.98, 1.34] & \textbf{2.72} [1.77, 4.73]\\
\bottomrule
\multicolumn{4}{@{}l}{\footnotesize $^\dagger$ conditional (suppressor) effect; not significant on its own (Section~\ref{sec:rq23}).}
\end{tabular}
\end{table}

\fig{fig_triangulation.png}{Triangulation of the five JDS skill areas across market demand, market
pay and junior salary hikes.}

Three messages follow. \textbf{Maths \& stats} is the only area that is both paid more in the market
and the strongest junior-hike signal. \textbf{Coding} is the baseline expectation (in 41\% of
data postings). \textbf{Dashboards \& storytelling} is rarely advertised and carries no clear market
premium, yet it is the second-strongest junior-hike signal: employers reward it after hiring more
than they ask for it in postings.

\textbf{Robustness check.} The result for storytelling depends on its definition. If ``reporting''
and ``MIS'' are counted as storytelling (the literal JDS wording), its market OR becomes 0.79
[0.69, 0.90], because those terms mark clerical MIS jobs. We report the narrow definition
(visualisation and BI tools, dashboards, storytelling) and disclose the alternative.
